% !TeX root = ../BTechProjectTemplate.tex
\chapter{Results \& Discussion}

\section{Quantitative Evaluation}
Swin-MMHCA performance will be assessed by PSNR and SSIM. The proposed method will be compared to a number of advanced approaches based on the data from the IXI database using T2 as the target modality.

\begin{table}[H]
\centering
\caption{Comparison of Super-Resolution Performance on IXI Dataset.}
\label{tab:quantitative_results}
\resizebox{\columnwidth}{!}{
\begin{tabular}{l|cc|cc}
\toprule
\textbf{Method} & \multicolumn{2}{c|}{\textbf{2$\times$}} & \multicolumn{2}{c}{\textbf{4$\times$}} \\
& PSNR $\uparrow$ & SSIM $\uparrow$ & PSNR $\uparrow$ & SSIM $\uparrow$ \\
\midrule
Bicubic & 33.44 & 0.9589 & 27.86 & 0.8611 \\
SRCNN \cite{srcnn} & 37.32 & 0.9796 & 29.69 & 0.9052 \\
VDSR \cite{vdsr} & 38.65 & 0.9836 & 30.79 & 0.9240 \\
IDN \cite{idn} & 39.09 & 0.9846 & 31.37 & 0.9312 \\
RDN \cite{rdn} & 38.75 & 0.9838 & 31.45 & 0.9324 \\
FSCWRN \cite{fscwrn} & 39.44 & 0.9855 & 31.71 & 0.9359 \\
CSN \cite{csn} & 39.71 & 0.9863 & 32.05 & 0.9413 \\
EDSR \cite{edsr} & 39.81 & 0.9865 & 31.83 & 0.9377 \\
\textbf{Swin-MMHCA (Ours)} & \textbf{40.62} & \textbf{0.9927} & \textbf{33.32} & \textbf{0.9572} \\
\bottomrule
\end{tabular}
}
\end{table}

Our Swin-MMHCA model achieves a superior PSNR and SSIM across both scales, demonstrating the success of the hierarchical transformer architecture in medical SR.

\section{Ablation Study and Modular Performance Analysis}
In order to assess the contribution of the different modules in the Swin-MMHCA architecture, an extensive ablation study was performed. We began with the base SwinV2 architecture and gradually incorporated the proposed modules.

\begin{itemize}
    \item \textbf{Baseline SwinV2 (Unimodal):} In the absence of any modality apart from the T2 images, the PSNR values were 38.12 dB for 2$\times$ upscaling and 31.14 dB for 4$\times$. The transformer had a good understanding of the image context, but not the high-frequency structural information that exists in other modalities.
    \item \textbf{Multi-modal Concatenation (Early Fusion):} Just by concatenating the T1 and PD modalities along with the T2 at the input, there was an improvement of 0.53 dB and 0.38 dB PSNR values for 2$\times$ and 4$\times$ upsampling respectively. The problem of reconciling the different distributions remained unsolved though.
    \item \textbf{MMHCA Module:} The addition of the Multi-Modal Hybrid Cross Attention module made the largest impact here, with values of 40.48 dB (+1.83 dB gain) for 2$\times$ and 32.90 DB (+1.38 dB gain) for 4$\times$. It helped in reconstructing sharp boundaries of the brain by querying relevant information from other scans.
    \item \textbf{Edge Guidance Branch:} The Sobel filter-based Edge guidance branch helped in further refinement of the structure boundaries, improving the values to 40.55 dB for 2$\times$ and 33.11 dB for 4$\times$. The Edge guidance branch served as a structure constraint mechanism.
    \item \textbf{Multi-Task Semantic Supervision:} However, adding the final stages of segmentation head and lesion detection head gave an improvement of 40.62 dB in 2$\times$ magnification and 33.32 dB in 4$\times$ magnification. More importantly, this step was effective in lowering the LPIPS score.
\end{itemize}

\section{Qualitative Analysis and Visual Results}
Visual inspection of the SR outputs across the test set confirms the quantitative findings. Swin-MMHCA demonstrates a superior ability to reconstruct the cerebral cortex's intricate sulci and gyri patterns. We provide visual evidence for both $2\times$ and $4\times$ upscaling factors.

\begin{figure}[htbp]
    \centering
    \includegraphics[width=\textwidth]{figures/comparison_figure_x4.png}
    \caption{Qualitative comparison of $4\times$ super-resolution results on the IXI T2-weighted brain MRI. From left to right: Low-Resolution (4x), Bicubic interpolation, EDSR-Nav (MHCA-main baseline), Swin-MMHCA (Ours), and High-Resolution Ground Truth. The zoom-in boxes highlight the recovery of sharp gray-white matter boundaries and fine anatomical textures.}
    \label{fig:results_visual_x4}
\end{figure}

\begin{figure}[htbp]
    \centering
    \includegraphics[width=\textwidth]{figures/comparison_figure_x2.png}
    \caption{Qualitative comparison of $2\times$ super-resolution results. At this lower upscaling factor, Swin-MMHCA achieves nearly perfect anatomical reconstruction, effectively bridging the gap between clinical-grade and high-resolution research scans.}
    \label{fig:results_visual_x2}
\end{figure}

As shown in Fig. \ref{fig:results_visual_x4}, Bicubic interpolation suffers from significant blurring, making it impossible to delineate the ventricular boundaries accurately. The EDSR-Nav model introduces "ringing" artifacts near high-contrast edges. In contrast, our Swin-MMHCA model produces a result that is visually indistinguishable from the ground truth for most subjects. The cross-attention mechanism successfully "borrows" the sharp contrast from the T1 modality to reconstruct the fine-grained details of the T2 scan, even at a high upscaling factor of $4\times$.

\section{Regional Performance Analysis}
To understand the model's strengths, we analyzed its performance across different anatomical regions of the brain.

\subsection{Cortical Sulci and Gyri}
The best enhancement was found in Swin-MMHCA compared to other baseline approaches, specifically in the cerebral cortex. It is important to note that the complex structures of the sulci are usually converted to a single gray block through bicubic interpolation. However, our model, utilizing the T1 structure as guidance, was able to accurately restore the boundary between the gray matter and cerebrospinal fluid.

\subsection{Ventricular Boundaries}
The lateral ventricles are high-contrast features that play an important role in monitoring brain atrophy during neurological disorders. The use of Swin-MMHCA ensured almost perfectly reconstructed lateral ventricles, with no issues of stair-cased artifacts associated with conventional CNNs.

\section{Clinical Implications of High-Fidelity SR}
The improvements obtained through Swin-MMHCA (such as attaining a PSNR of 40.62 dB) have significant implications in the medical field. Feature extraction in radiomics is highly sensitive to noise and resolution in the image. Through generating super-resolution images that maintain consistency in structure with ground truth high-resolution images, our framework allows for:
\begin{enumerate}
    \item \textbf{More Accurate Volumetry:} Reducing the partial volume effect leads to more precise measurements of hippocampal volume, a key biomarker for Alzheimer's disease.
    \item \textbf{Improved Lesion Detection:} Sharp textures allow for the detection of subtle micro-lesions in Multiple Sclerosis that might be missed in low-resolution clinical scans.
    \item \textbf{Surgical Planning:} High-resolution maps of the cortical surface are essential for neurosurgical navigation and avoiding critical functional areas.
\end{enumerate}
Our results suggest that Swin-MMHCA can effectively bridge the gap between "fast clinical protocols" and "high-resolution research standards," potentially reducing the need for long, expensive scans.
